\documentclass[11pt]{article}
\usepackage{amsmath,amssymb}
\usepackage{graphicx}
\usepackage{hyperref}

\title{Language as Telemetry: Wittgenstein, Structural Signals, and Admissible Meaning}
\author{J. M. Bookbinder \\ Witness Grade Analytics \\ Atlanta, GA, USA}
\date{2026}

\begin{document}
\maketitle

\begin{abstract}
Operational systems increasingly rely on language to initiate, control, and document state transitions. 
This paper argues that such language must be treated not as narrative but as telemetry: a structured signal 
whose admissibility depends on source integrity, rule-governed transformation, and replay-verifiable witness preservation.
Drawing on Wittgenstein’s account of rule-following and language games, we show that meaning in operational contexts 
is not a private interpretation but a public, rule-governed activity. 
This perspective aligns naturally with continuation-control architectures in which system execution proceeds only when 
admissibility conditions are satisfied. 
\end{abstract}

\section{Introduction}

Modern software systems frequently treat natural language as an operational interface. 
Prompts, commands, and policy text may initiate workflows, trigger state changes, or alter system behavior. 
When language operates in this capacity it functions not merely as description but as signal.

A signal transmitted through an operational system must satisfy structural constraints:
source identity, transformation rules, temporal ordering, and replay-verifiable preservation. 
In engineering terms, this is telemetry.

This paper proposes that language used within operational systems must therefore be treated as telemetry rather than narrative.

\section{Wittgenstein and Rule-Governed Meaning}

Wittgenstein argued that meaning arises from rule-governed use within a language game. 
Understanding a statement is not a private mental act but participation in a public system of rules.

In operational environments this insight has a direct analogue. 
A statement that triggers execution must be interpreted according to explicit rules and system state. 
Meaning is therefore constrained by admissibility conditions rather than interpretive freedom.

\section{Telemetry as a Model of Language}

Telemetry systems transmit structured signals from instruments to receivers. 
Each signal possesses:

\begin{itemize}
\item A source address
\item A channel
\item A timestamp
\item A transformation history
\end{itemize}

Language used as an operational command possesses the same structure:

\begin{itemize}
\item An author or origin
\item A communication channel
\item A temporal ordering
\item A transformation path through the system
\end{itemize}

This structural similarity motivates treating language as telemetry.

\section{Admissibility and Continuation Control}

Continuation-control architectures determine whether a declared process may proceed under explicit constraints.

Let a continuation $c$ be admissible when all required structural conditions hold.

\[
\mathrm{Adm}(c) = 1 \iff
\begin{cases}
\text{time reconciled} \\
\text{drift bounded or paid} \\
\text{gaps preserved and owned} \\
\text{operator binding present} \\
\text{deterministic replay holds}
\end{cases}
\]

Language signals that trigger execution must therefore satisfy the same admissibility conditions.

\section{Witness Preservation}

To maintain admissibility, operational language must produce verifiable artifacts:

\begin{itemize}
\item canonicalized representations
\item cryptographic hashes
\item deterministic replay instructions
\item refusal records when admissibility fails
\end{itemize}

These artifacts serve as witnesses certifying that meaning has remained intact under transformation.

\section{Conclusion}

Language used in operational systems is best understood as telemetry: a rule-governed signal whose admissibility 
depends on structural verification rather than narrative interpretation. 
By treating language as telemetry, systems can enforce continuation-control architectures that ensure execution proceeds 
only under replay-verifiable evidence.

\end{document}
